\documentclass[12pt,a4paper]{article}

\usepackage[top=2cm, bottom=2cm, left=2cm, right=2cm]{geometry}
\usepackage{fontspec}
\setmainfont{Times New Roman}
\usepackage{xeCJK}
\setCJKmainfont{PingFang TC}
\usepackage{xcolor}
\usepackage{booktabs}
\usepackage{longtable}
\usepackage{amssymb,amsmath}
\usepackage{enumitem}
\usepackage{hyperref}

\definecolor{errorred}{RGB}{200,0,0}
\definecolor{warncolor}{RGB}{180,100,0}
\definecolor{okgreen}{RGB}{0,128,0}
\definecolor{resolved}{RGB}{0,100,180}

\hypersetup{colorlinks=true, linkcolor=blue, citecolor=blue, urlcolor=blue}

\begin{document}

\begin{center}
{\Large\bfseries Academic Review Report R3}\\[0.5em]
{\large Volatility Absorption: The Diminishing Marginal Impact of Market Fear}\\[0.5em]
{\normalsize Document type: Empirical finance paper (target: JBF / JFQA)}\\
{\normalsize File: \texttt{paper/volatility-absorption/main\_v2.tex}}\\
{\normalsize Review date: 2026-04-05}\\
{\normalsize Review standard: Strict (journal-referee level)}\\
{\normalsize Reference: R2 review dated 2026-04-05; K897 null simulation; K903 robustness attempt (no results JSON)}
\end{center}

\bigskip

%% ============================================================
\section*{R3 Review Summary}

\begin{tabular}{ll}
\textcolor{okgreen}{R1 SEVERE issues resolved} & 2 of 5 (S1, S5) \\
\textcolor{resolved}{R1 SEVERE downgraded} & 1 of 5 (S3 $\to$ MEDIUM) \\
\textcolor{errorred}{R1 SEVERE issues still open} & 2 of 5 (S2, S4) \\
\textcolor{warncolor}{New issues identified in R3} & 1 item (N2: 2020--2026 sign reversal) \\
Overall Assessment & \textbf{Major Revision (narrowed scope)} \\
\end{tabular}

\bigskip

%% ============================================================
\section*{R3 Scorecard}

\begin{tabular}{lll}
\toprule
\textbf{Issue} & \textbf{R3 Status} & \textbf{Severity} \\
\midrule
S1: No null simulation & \textcolor{okgreen}{RESOLVED (R2)} & --- \\
S2: Table 5 sample-size inconsistency & \textcolor{errorred}{STILL OPEN --- data vintage} & SEVERE \\
S3: Tables 9--10 untraceable & \textcolor{resolved}{DOWNGRADED --- content present, partially reproduced} & MEDIUM \\
S4: NFP table discrepancies with K741 & \textcolor{errorred}{STILL OPEN} & SEVERE \\
S5: Orphan reference \texttt{chernov2018} & \textcolor{okgreen}{RESOLVED (R2)} & --- \\
\addlinespace
N1: Robustness tables lack experiment JSON & \textcolor{warncolor}{UNCHANGED} & MEDIUM \\
N2 (NEW): 2020--2026 sub-period sign reversal & \textcolor{errorred}{NEW} & MAJOR \\
\bottomrule
\end{tabular}

\medskip

\noindent\textbf{SEVERE count: 2} (down from 3 in R2, 5 in R1).\\
\textbf{Overall assessment: Major Revision (narrowed scope).} The core SAR finding is robust (K897 decisively rejects the GARCH null, $p = 0.002$). The remaining severe issues are data-integrity problems (sample sizes, NFP reconciliation) rather than methodological threats. However, the newly identified 2020--2026 sign reversal in the sub-period robustness raises a genuine scientific concern that requires honest treatment.

\bigskip

%% ============================================================
\section*{Detailed Assessment of Each Issue}

\subsection*{S2: Table 5 Sample-Size Inconsistency \hfill \textcolor{errorred}{SEVERE}}

\textbf{R2 status:} Table~5 reports $N = 127$ (rate shocks), $N = 203$ (risk-off), $N = 89$ (geopolitical). K721 stores different values ($79$, $182$, $146$). None match.

\textbf{R3 assessment:} The robustness verification attempt (K903 plan) determined that the $N$ values are \textbf{unreproducible due to data vintage change}. When the analysis is re-run with current yfinance data, the number of shock days falling into each type category shifts because:
\begin{enumerate}[nosep]
\item VIX and price data undergo minor revisions (corporate actions, splits adjustments).
\item The sample end date has extended, adding new shock-day classifications.
\item The boundary between VIX regimes is sensitive to the exact VIX values on marginal days.
\end{enumerate}

This is a known problem in empirical finance: rerunning an analysis months later with fresh data downloads can produce slightly different sample compositions. However, it does \textbf{not} excuse the discrepancy between the paper's $N$ values and K721 (which was the original experiment). The paper either used (a)~a different computation than K721, or (b)~a corrected version of K721 that was not saved.

\textbf{Why this remains SEVERE:} The shock-type decomposition is ``the paper's most important result'' (Section~5.4, line 397 of \texttt{main\_v2.tex}). The central claim---endogenous shocks are absorbed, exogenous shocks are not---rests entirely on this table. If the $N$ values, absorption coefficients, and $t$-statistics cannot be traced to a specific reproducible computation, the paper's core contribution is empirically unanchored.

\textbf{Required action:}
\begin{enumerate}[nosep]
\item Create a new experiment (e.g., \texttt{k9XX\_absorption\_shock\_types.py}) that computes Table~5 from scratch with current data.
\item Document the exact shock-classification methodology: how are rate shocks, risk-off flights, and geopolitical shocks defined? What are the sign thresholds for SPY, TLT, GLD? Is $|\Delta\text{VIX}| > 2$ the only filter, or is there also a $r_t^{\text{SPY}} < 0$ requirement?
\item Report the new $N$ values in the paper. If they differ from the current table, update the table.
\item Store the bootstrap $t$-statistics and absorption coefficients in the experiment JSON.
\item If the original K721 values and the current values differ substantively (i.e., the absorption coefficients change sign or lose significance), this must be reported honestly.
\end{enumerate}


\subsection*{S3: Tables 9--10 Robustness \hfill \textcolor{resolved}{DOWNGRADED to MEDIUM}}

\textbf{R2 status:} SEVERE---no experiment JSON, untraceable numbers.

\textbf{R3 assessment:} The paper now contains substantive robustness tables:
\begin{itemize}[nosep]
\item Table~\ref{tab:robust_threshold}: 5 alternative shock thresholds ($\tau = 1, 1.5, 2, 2.5, 3$). All negative and significant ($p < 0.03$ for all). Monotonic increase in $|\hat{\beta}|$ with $\tau$---internally consistent.
\item Table~\ref{tab:robust_subperiod}: 3 sub-periods. All negative. GFC strongest ($-0.00035$), low-vol weakest ($-0.00018$, marginally significant at $p = 0.087$), COVID era significant ($-0.00031$, $p = 0.008$).
\item In-text: RV normalization ($\hat{\beta}^{\text{RV}} = -0.0031$, $t = -2.76$) and controlled regression ($\hat{\beta} = -0.00025$, $t = -3.14$).
\end{itemize}

The robustness attempt (K903) \textbf{partially reproduced} these results:
\begin{itemize}[nosep]
\item Threshold variations: Reproduced with consistent signs, though with \textbf{weaker $t$-statistics} than the paper reports. This is expected if the data vintage changed slightly.
\item Sub-period stability: 2006--2012 and 2013--2019 reproduced directionally. \textbf{However, 2020--2026 showed a sign reversal} (see N2 below).
\end{itemize}

\textbf{Downgrade rationale:} The robustness content is now present in the paper and the \textit{overall pattern} is largely confirmed. The threshold monotonicity and the full-sample significance are robust. The absence of a backing experiment JSON is a process issue (N1, MEDIUM), not a scientific one. The sub-period sign reversal is elevated as a new issue (N2, MAJOR) rather than keeping S3 at SEVERE.


\subsection*{S4: NFP Table Discrepancies with K741 \hfill \textcolor{errorred}{SEVERE}}

\textbf{R2 status:} SEVERE---paper's Table~6 numbers ($n = 63$, overall $p = 0.037$) do not match K741 JSON ($n = 62$, $p = 0.081$ or $0.061$).

\textbf{R3 assessment:} No reconciliation has been performed. The K903 robustness attempt did not address NFP. The discrepancies remain:

\medskip
\begin{tabular}{llrrl}
\toprule
\textbf{Value} & \textbf{Paper} & \textbf{K741} & \textbf{Difference} & \textbf{Concern} \\
\midrule
Low VIX $n$ & 63 & 62 & $+1$ & Minor \\
Overall $p$-value & 0.037 & 0.081 or 0.061 & $-0.024$ to $-0.044$ & \textcolor{errorred}{Material} \\
Elevated $|r|$ & 1.053 & 1.022 & $+0.031$ & Moderate \\
\bottomrule
\end{tabular}

\medskip

The overall $p$-value discrepancy is the most concerning. The paper claims significance at 5\% ($p = 0.037$), while K741 suggests $p \approx 0.06$--$0.08$ (not significant at 5\%). If the correct value is $p > 0.05$, the paper overstates the evidence from the NFP analysis.

The paper's footnote (line 70) specifying ``195 NFP release dates over 4,081 SPY trading days'' is a helpful addition, and the limitation paragraph about surprise controls is commendable. But these do not resolve the numerical discrepancy.

\textbf{Required action:}
\begin{enumerate}[nosep]
\item Create a new experiment that computes the NFP analysis from scratch.
\item Use the BLS release calendar to identify NFP dates (the paper describes this methodology).
\item Report the new values. If $p > 0.05$, the paper must weaken its NFP claim accordingly (``directionally consistent but not statistically significant at conventional levels'').
\item Note that the NFP analysis is supplementary to the SAR-based primary finding; weakening the NFP claim does not undermine the core contribution.
\end{enumerate}


\subsection*{N2 (NEW): 2020--2026 Sub-Period Sign Reversal \hfill \textcolor{errorred}{MAJOR}}

\textbf{Finding:} The K903 robustness attempt found that when the sub-period analysis is re-run with current data, the \textbf{2020--2026 sub-period shows a sign reversal}---the absorption coefficient becomes positive or near-zero, contradicting the paper's Table~\ref{tab:robust_subperiod} which reports $\hat{\beta} = -0.00031$ ($t = -2.65$, $p = 0.008$).

\textbf{Why this is concerning:}
\begin{enumerate}[nosep]
\item The paper presents 2020--2026 as \textit{supporting evidence} for the robustness of the absorption effect across sub-periods. If the sign actually reverses, this is false.
\item The 2020--2026 period is the most recent and arguably most policy-relevant. If absorption does not hold in the post-COVID era, this significantly limits the paper's practical implications (Section~6: hedging cost-benefit, VT strategy design).
\item The sign reversal may be due to the unusual VIX dynamics of 2020--2026: the extreme COVID spike (VIX $> 80$) followed by an unprecedented rapid normalization, then the 2022 rate-hiking regime with elevated but not extreme VIX. These are structurally different from the GFC era.
\end{enumerate}

\textbf{Possible explanations (to investigate):}
\begin{itemize}[nosep]
\item \textbf{Data vintage:} The paper's computation used data downloaded in early 2026; the K903 attempt used a later download. If a few marginal shock days shifted classification, the sub-period result (with fewer observations) could flip.
\item \textbf{Sample composition:} The 2020--2026 period contains both the extreme COVID outlier (which may dominate the regression) and the 2022--2023 hiking cycle. Excluding March 2020 or splitting 2020--2021 from 2022--2026 could reveal whether the reversal is driven by a single outlier month.
\item \textbf{Genuine structural change:} The post-COVID market may have experienced a genuine shift in absorption dynamics---perhaps because the speed of the COVID recovery (V-shaped) was unlike previous crises, or because the 2022 rate-hiking regime produced a different shock structure.
\end{itemize}

\textbf{Required action:}
\begin{enumerate}[nosep]
\item Create a formal experiment re-running the sub-period analysis with current data.
\item If the 2020--2026 sign reversal is confirmed:
\begin{itemize}[nosep]
\item Report it honestly in the paper. The robustness section must include the reversal.
\item Test whether the reversal is driven by the March 2020 outlier (re-run excluding 2020-02 to 2020-04).
\item Split 2020--2026 into 2020--2022 and 2023--2026 to isolate the source.
\item Discuss possible structural explanations (COVID V-shape recovery, rate-hiking regime).
\end{itemize}
\item If the reversal is a data-vintage artifact, document this and use a \textbf{pinned dataset} (specific download date, saved locally) for all paper computations.
\item \textbf{Do not suppress the finding.} If the absorption effect weakened or reversed in the most recent period, this is a genuine empirical observation that deserves discussion, not concealment. The paper's core finding (full-sample SAR, K897 null rejection) remains valid even if one sub-period is anomalous.
\end{enumerate}


\bigskip

%% ============================================================
\section*{Status of R1/R2 MEDIUM Issues}

\begin{longtable}{p{1cm}p{5cm}p{2.5cm}p{5.5cm}}
\toprule
\textbf{ID} & \textbf{Issue} & \textbf{R3 Status} & \textbf{Assessment} \\
\midrule
\endfirsthead
M1 & Novelty positioning (MS-GARCH) & \textcolor{warncolor}{UNCHANGED} & K897 shows standard GARCH cannot produce SAR decline. MS-GARCH remains untested. \\
\addlinespace
M2 & Missing literatures (VVIX, HAR-RV) & \textcolor{warncolor}{UNCHANGED} & \\
\addlinespace
M3 & Danielsson et al.\ (2018) citation & \textcolor{okgreen}{IMPROVED (R2)} & Adequate. \\
\addlinespace
M4 & Table 1 (descriptive stats) & \textcolor{warncolor}{UNCHANGED} & No replication script. \\
\addlinespace
M5 & 767 vs.\ 893 sample-size discrepancy & \textcolor{warncolor}{UNCHANGED} & \\
\addlinespace
M6 & \texttt{lin2020} unverifiable & \textcolor{warncolor}{UNCHANGED} & \\
\addlinespace
M7 & \texttt{bekaert2022} end page & \textcolor{warncolor}{UNCHANGED} & \\
\addlinespace
M8 & Replication scripts for K716--K721 & \textcolor{warncolor}{UNCHANGED} & \\
\addlinespace
N1 & Robustness tables lack JSON & \textcolor{warncolor}{UNCHANGED} & Create \texttt{k9XX\_absorption\_robustness.json} \\
\bottomrule
\end{longtable}


\bigskip

%% ============================================================
\section*{Honest Assessment: What Can and Cannot Be Downgraded}

\subsection*{S3: Can be downgraded (SEVERE $\to$ MEDIUM). Done.}

Rationale: The robustness content is substantive. The threshold monotonicity is confirmed. The full-sample result ($\hat{\beta} = -0.00028$, $t = -3.42$) is robust. The process gap (no JSON) is a MEDIUM concern. The sign reversal is elevated as a separate issue (N2).

\subsection*{S2: Cannot be downgraded. Remains SEVERE.}

Rationale: The shock-type decomposition is the paper's ``most important result'' (author's own words, line 397). If the $N$ values, absorption coefficients, and $t$-statistics cannot be independently verified, the core contribution is empirically unanchored. The data-vintage explanation is plausible but does not resolve the discrepancy---it merely explains \textit{why} the numbers might differ, not which numbers are correct. A fresh experiment with documented methodology is required.

\subsection*{S4: Cannot be downgraded. Remains SEVERE.}

Rationale: The overall $p$-value discrepancy ($0.037$ paper vs.\ $0.06$--$0.08$ K741) crosses the conventional significance threshold. If the true $p > 0.05$, the paper misrepresents the evidence. This is a binary issue---significant or not---that cannot be finessed.

\subsection*{N2: Elevated to MAJOR (not SEVERE).}

Rationale: The 2020--2026 sign reversal is concerning but does not invalidate the core finding. The full-sample SAR analysis and K897 null simulation are unaffected. The sign reversal in one sub-period suggests the absorption effect may be time-varying---an important nuance that strengthens rather than weakens the paper \textit{if reported honestly}. It would be SEVERE only if the paper claimed unconditional stability and refused to acknowledge the reversal.


\bigskip

%% ============================================================
\section*{R3 Revision Priority}

\textbf{Priority 1 (must fix before submission):}
\begin{enumerate}
\item \textbf{Re-run Table~5 shock-type analysis} with current data, documented methodology, and experiment JSON. Update $N$ values and coefficients if they change. \textit{[S2]}
\item \textbf{Re-run NFP analysis} with documented BLS calendar dates. If $p > 0.05$, weaken the claim. \textit{[S4]}
\item \textbf{Investigate and report the 2020--2026 sign reversal.} Run formal sub-period tests, test outlier sensitivity (exclude March 2020), and discuss structural explanations. \textit{[N2]}
\end{enumerate}

\textbf{Priority 2 (strongly recommended):}
\begin{enumerate}[resume]
\item \textbf{Create unified robustness experiment JSON} containing all robustness results (thresholds, sub-periods, RV normalization, controlled regression). \textit{[S3, N1]}
\item \textbf{Pin the dataset:} Save a local copy of the exact data used for all paper computations. Document the download date and yfinance version. This prevents future data-vintage discrepancies. \textit{[S2, S4 prevention]}
\item \textbf{Create replication scripts} for K716--K721. \textit{[M8]}
\end{enumerate}

\textbf{Priority 3 (desirable):}
\begin{enumerate}[resume]
\item Add VVIX and HAR-RV literature discussion. \textit{[M2]}
\item Verify \texttt{lin2020} citation. \textit{[M6]}
\item Explain 767 vs.\ 893 sample-size gap. \textit{[M5]}
\item Test MS-GARCH null. \textit{[M1]}
\item Present RV normalization and controlled-regression results in a table. \textit{[N1]}
\end{enumerate}


\bigskip

%% ============================================================
\section*{What the Paper Does Well (Unchanged)}

The core strengths from R2 remain:
\begin{enumerate}
\item \textbf{K897 null simulation is exemplary.} Cohen's $d > 3.8$ in all bins, $p = 0.002$. The SAR finding is not a GARCH artifact.
\item \textbf{SAR measure is clean.} No VIX denominator, no mechanical correlation.
\item \textbf{Endogenous/exogenous distinction is novel.} If the Table~5 numbers hold up to re-verification, this is a genuine contribution.
\item \textbf{NFP limitation paragraph is honest.} The surprise-control caveat is commendable.
\item \textbf{Robustness section is substantive.} Five thresholds, three sub-periods, two alternative normalizations.
\end{enumerate}


\bigskip

%% ============================================================
\section*{Overall Assessment}

\textbf{Verdict: Major Revision (narrowed scope).}

The paper has made critical progress since R1 (5 SEVERE $\to$ 2 SEVERE). The methodology is sound (K897 definitively rejects the mechanical null). The remaining problems are \textbf{data integrity} (S2: shock-type $N$ values, S4: NFP $p$-value) and \textbf{honest reporting} (N2: 2020--2026 sign reversal). None of these threaten the core SAR finding or the null-simulation identification.

The path to submission is clear:
\begin{enumerate}[nosep]
\item Re-run Tables 5 and 6 with documented methodology $\to$ update numbers $\to$ if significance changes, adjust claims.
\item Investigate and honestly report the 2020--2026 sub-period finding $\to$ either confirm it was a data-vintage artifact or discuss it as time-variation in absorption.
\item Pin the dataset to prevent future discrepancies.
\end{enumerate}

With these three actions, the paper would reach SEVERE = 0 and be submittable to JBF/JFQA. The core contribution---SAR, endogenous/exogenous decomposition, and null-simulation identification---is publishable.

\bigskip

\noindent\textit{Reviewer: Claude Opus 4.6 (1M context), acting as JBF/JFQA referee. R3 review, 2026-04-05.}\\
\noindent\textit{R2 reference: \texttt{paper/volatility-absorption/reviews/review\_r2.tex}}\\
\noindent\textit{K897 source: \texttt{paper/volatility-absorption/experiments/k897\_sar\_null\_simulation\_results.json}}\\
\noindent\textit{K903 robustness attempt: results JSON not found; findings based on AUDIT\_PLAN.md and user report.}

\end{document}
